\section{Octava iteración de la Fase 5: Las tres imágenes y el GWC solo para ejecución}
\label{sec:F5I8}

Esta iteración cierra dos frentes abiertos por el uso real del
simulador: el contrato de las imágenes de pantalla en la validación
de entrada, que quedaba a medias tras hacer opcional la imagen de
fondo, y la carga del GWC, que jugaba pero dejaba la máquina a
medias (botones sin restaurar y pestañas heredadas del cartucho
anterior).

\subsection{El contrato de las tres imágenes}
\label{sec:F5I8Imagenes}

La validación de entrada (\texttt{validarCartucho}, en
\texttt{BaseCartucho.js}) es el contrato de lo que puede ENTRAR, y es
COMÚN a los cartucho JSON: GWB y GWI (el GWC es binario y se valida
por su propio camino, en la lectura del fichero). Con ella se fija el
contrato de las tres imágenes de pantalla, que el anexo
\ref{sec:Formatos} recoge ya de forma formal:

\begin{itemize}
	\item \texttt{pantalla} es OBLIGATORIA como objeto: sin ella no
	hay ni geometría ni presentación.
	\item \texttt{pantalla.img\_s} es OBLIGATORIA: de ella se extrae
	la geometría, así que debe estar declarada Y su fichero presente
	en la carpeta o en el ZIP. Un GWI cuyo ZIP no la lleve ya no
	carga.
	\item \texttt{pantalla.img\_f} (el fondo) es OPCIONAL (punto 8 de
	la Fase 5): hay cartuchos sin arte de fondo. Si se declara, su
	fichero DEBE estar presente; sin fondo, el simulador pinta la
	pantalla en blanco (el color real del LCD apagado).
	\item \texttt{pantalla.img\_m} (la portada) es OPCIONAL, y AHORA
	TAMBIÉN se valida: si se declara, su fichero debe estar. Antes una
	portada ausente se toleraba en silencio al cargar.
\end{itemize}

La regla común es simple y uniforme: una imagen es opcional a la
hora de DECLARARSE, pero declarada su fichero es obligatorio. El
guardado del GWI se alinea con el mismo criterio: el ZIP incluye
img\_f solo si el cartucho la declara; el aviso de error se reserva
al caso incoherente (declarada y sin imagen en memoria), y un
cartucho sin fondo deja solo constancia en consola, porque es
legítimo.

\subsection{El GWC como cartucho de pura ejecución}
\label{sec:F5I8GWC}

Cargar el GWC generado y comprobarlo en el PC es parte del ciclo de
trabajo: es jugable y sirve para verificar que el binario es
correcto. Pero el GWC NO lleva la parte PC del cartucho: no trae
img\_s (sin geometría no hay mapa de trabajo ni autoría) ni los
botones con sus nombres. La carga quedó por tanto definida como
EJECUCIÓN PURA, con tres ajustes en \texttt{procesarGWC}:

\begin{enumerate}
	\item \textbf{Botonera completa}: el GWC no declara botones (son
	físicos del setup), así que se restaura la botonera COMPLETA de la
	máquina, por si un cartucho anterior había eliminado o renombrado
	alguna, y el botón de Start vuelve al valor por defecto (el 8).

	\item \textbf{Pestañas}: con un GWC cargado, solo la pestaña del
	Simulador está activa; Editor y Segmentos y Sprites quedan
	deshabilitados, porque no hay nada de la parte PC que editar ni
	listar de forma fiable.

	\item \textbf{La máquina delante}: al terminar la carga, la vista
	Simulador queda delante y el juego listo para el Start, venga de
	la pestaña que viniera.
\end{enumerate}

Se valoró reconstruir el estado del editor a partir de las listas
del GWC (sprites desde sus máscaras, segmentos con sus propiedades,
historial vacío) para poder editar un GWC sin GWI al lado; y también
un fichero secundario con la parte PC. Ambos caminos se descartaron
por innecesarios: el par (GWI, GWC) que el botón «Generar GWI+GWC»
ya produce ES la pareja de parte PC y parte de ejecución, y el GWC
solo cumple su papel: verificar el binario jugando. La edición de un
GWC sin su GWI queda anotada como AMPLIACIÓN del proyecto.